\chapter{Data and Training Setup}
\label{ch:data_setup}
This section specified the dataset characteristics and training choices. A summary table is made available at Tab.\ref{tab:training_setup}. 
\section{Data}

The dataset follows the source and structure reported by Kim et al.\ \cite{kim2025diffusion}.
Specifically, the list of current S\&P 500 members was taken as the starting universe (as of April 2026), and
tickers with non-standard symbols (e.g., BRK.B, BF.B) were excluded. For each remaining
ticker, the maximum available daily price history was downloaded. The dataset was then
filtered to retain only those stocks whose available history extends at leas 40 years back.
According to the authors:
\begin{quote}
    \emph{This selection criterion ensures that the dataset encompasses diverse market
    regimes, including various boom and bust cycles, thus providing a rich context for
    training the generative model.}
\end{quote}
For each selected stock, daily log-returns were computed as in Eq.~\ref{eq:log_returns},
using the adjusted closing price on day \(t\).

This selection is subject to some limitations. First, survivorship bias: only companies that remained listed for a
period substantially longer than the average tenure in the S\&P 500 index, which is 15 years \cite{slok_apollo_longevity}, are included. The resulting universe therefore may understate the frequency of adverse market events associated with delistings
and corporate failures. The second aspect regards the decimalization process, which was mandated by the Security and Exchange Commission in January 2000 and completed by the 9th of April 2001 \cite{furfine2003decimalization}. Before this period US stock prices were quoted in fractions, specifically 1/16th of a dollar, meaning that the minimum price movement needed to record an effective change in price was 0.0625 dollars. After decimalization this changed to 1 cent. This manifest as 7.1928\% of all returns displaying an exact 0\% daily returns.

\section{Preprocessing}

In Phase 1, only univariate log-returns are used, as defined in Eq.~\ref{eq:log_returns}.
In Phase 2, where conditioning on all four OHLC channels is required, the Open, High,
and Low prices are each expressed as log-returns relative to the previous closing price:
\[
    \hat{O}_t = \log\!\left(\frac{O_t}{C_{t-1}}\right), \qquad
    \hat{H}_t = \log\!\left(\frac{H_t}{C_{t-1}}\right), \qquad
    \hat{L}_t = \log\!\left(\frac{L_t}{C_{t-1}}\right)
\]
This formulation is consistent with the Close log-return defined in Eq.~\ref{eq:log_returns},
expressing all channels relative to the same reference price \(C_{t-1}\).

\(\textbf{Phase~2}\quad\) A per-channel z-score normalisation is then applied independently to each of the four
distribution while standardising it to zero mean and unit variance, so that
\(\sigma_{\text{data}} = 1\) holds by construction and network inputs remain numerically
stable across channels. However, because each channel is scaled by its own standard
deviation, the ordering constraint \(L_t \leq \min(O_t, C_t) \leq \max(O_t, C_t) \leq H_t\),
which holds in price space, is not preserved after normalisation. This is a structural
limitation acknowledged throughout the evaluation.

Four tickers were excluded from the dataset prior to training: BEN, HUBB, NVR, 
and WMB. Each exhibits absolute log-returns exceeding $100\%$, with values 
ranging from $-157\%$ to $+330\%$. Returns below $-100\%$ are theoretically 
inadmissible, because they violate price non-negativity. Such observations are therefore artifacts of the split- and dividend-adjustment procedure described in 
Sec.~\ref{sec:fts_literature}, rather than genuine market events. Exclusion was 
preferred over imputation or smoothing, both of which would introduce their own 
distortions. This choice is consistent with the likely preprocessing of Kim 
et al.~\cite{kim2025diffusion}: after removing these four tickers, the empirical 
distribution aligns more closely with Figure~8 of that paper, 
which is used here as a reference baseline.

\section{Sliding Window Construction}
\label{sec:sliding_window_ch4}

Training sequences are constructed via a sliding window over the full dataset. Two
configurations are used across the two experimental phases:
\begin{itemize}
    \item \textbf{Phase 1, replication:} window length \(L = 2048\) and stride \(s = 400\),
    yielding approximately 5,500 training windows over the full dataset.
    \item \textbf{Phase 2, EDM-CSDI:} window length \(L = 512\) and stride \(s = 100\), yielding approximately 23,000
    training windows cosistent with what described in Ch.~\ref{ch:theoretical_framework}.
\end{itemize}

A window of length \(L = 2048\) spans approximately ten years of daily returns and is
therefore well-suited to capturing long-term trends. However, the temporal Transformer
incurs an \(\mathcal{O}(L^2)\) cost in sequence length, making it computationally
expensive. The shorter \(L = 512\) setting was adopted for Phase 2 due to hardware
constraints imposed by MIG partitioning of the A100 GPU by the High--Performance Computing's university policy. The full memory analysis is given in App.~\ref{app:sec:attention_memory}.

\(\textbf{Phase 1} \quad\) Given the window length the number of non-overlapping windows per ticker is small, making a strict 
train--validation split undesirable. Validation windows were therefore extracted 
by random holdout, consistent with the supposed implementation of Kim 
et al.~\cite{kim2025diffusion}, who are assumed to have either applied a similar random selection or foregone a validation set entirely. Due to the stride, partial overlap between training and validation windows exists.

For the GBM-inspired SDE variant, each window was additionally mean-centred prior 
to training by subtracting the sample mean of the window. This transformation does not conflict with the generative setting, because no forecasting is performed. In addition, any daily log-returns can be easily recovered by \(r_{\tau+k}=C_{\tau+k}-C_{\tau+k-1}\), except for the first element.

\(\textbf{Phase 2}\ \quad \) In order to assess the discriminative power of the conditioning information, validation windows were fully held out: for any ticker whose window was selected for validation, no other window covering the same time range appeared in the training set. This reduces the number of training samples available for learning, however, this effect is mitigated by the larger number of windows available.

% Overlapping windows introduce correlation between adjacent training samples, violating
% the i.i.d.\ batch assumption. This is a standard and accepted approximation in
% deep learning on time series, analogous to patch sampling in image models, but it should
% be noted as a limitation when interpreting training dynamics.

[INTRODUCE WINDOW'S CREATION DIFFERENCE AND include this also in the TABLE for set up; include also sampler for fair comparison]

\section{Training Configuration}
\label{sec:training_config}

This section summarises the most relevant training configuration and design choices;
explicit formulations for the optimiser and learning rate schedule are given in
App.~\ref{app:sec:optim}, and full hyperparameter details in App.~\ref{app:sec:training_config}.

\textbf{Optimizer} \quad The optimiser used throughout is \texttt{AdamW}~\cite{loshchilov2019decoupled},
following the choice of CSDI~\cite{tashiro2021csdi}. It applies weight decay directly
to the parameters rather than through the gradient, preventing the interaction between
$\ell_2$ regularisation and the adaptive scaling of standard Adam. This decoupling results in more consistent regularisation across parameters with different
gradient magnitudes.

\textbf{Learning Rate Scheduler} \quad The learning rate follows cosine
annealing~\cite{loshchilov2017sgdr}, allowing larger updates early in training and progressively smaller updates near
convergence. Reduce-on-Plateau was also evaluated but yielded consistently worse
training dynamics and was not adopted. For Phase~1, $\eta_{\max} = 5 \times 10^{-4}$
and $\eta_{\min} = 1 \times 10^{-6}$; for Phase~2, $\eta_{\max}$ is lowered to
$1 \times 10^{-4}$ to reflect the smaller effective gradient variance under EDM
loss weighting.

\textbf{Model Architecture} \quad  Both phases share the same backbone: 128 \texttt{channels},
6 \texttt{layers}, a 256-dimensional diffusion-step and time embedding, and a
64-dimensional feature embedding, following Kim et al.~\cite{kim2025diffusion}.
This configuration was selected empirically: a 64-channel variant was less expressive,
while a 256-channel variant offered no measurable improvement in generation quality at
substantially higher compute cost.

\textbf{Noise Schedule} \quad The noise hyperparameters differ between phases.
In Phase~1, they define the range of the forward-process noise: VE-style diffusion
uses $\sigma_{\max} = 1$ and $\sigma_{\min} = 0.01$; GBM-style diffusion, operating
on cumulative log-returns, uses $\sigma_{\max} = 10$ and $\sigma_{\min} = 0.1$; and
the VP formulation is parameterised by $\beta_{\max} = 8$ and $\beta_{\min} = 0.01$.
In Phase~2, the EDM log-normal sampling distribution
$\ln\sigma \sim \mathcal{N}(P_{\text{mean}}, P_{\text{std}}^2)$ replaces the
forward-process schedule entirely. The values $P_{\text{mean}} = -1.4$ and
$P_{\text{std}} = 1.6$ were selected via grid search over
$P_{\text{mean}} \in \{-0.8,\,-1.0,\,-1.2,\,-1.4,\,-1.6\}$ and
$P_{\text{std}} \in \{1.0,\,1.2,\,1.4,\,1.6,\,1.8\}$ based on validation loss;
the full ablation is reported in Sec.~\ref{sec:phase2_ablation}.
\textbf{Training parameters.} The batch size is 64 across both phases. Phase~1 trains for the same number of epochs as Kim et al.~\cite{kim2025diffusion}, while Phase~2 uses half that budget, which reflects the hardware constraints and the larger number of available training windows, which increases the number of gradient steps per epoch.

\textbf{Sampling parameters} \quad For Phase~2, $\sigma_{\max}$ and $\sigma_{\min}$ were
selected to cover the support of $\ln\sigma \sim \mathcal{N}(P_{\text{mean}}, P_{\text{std}}^2)$,
requiring $\sigma_{\max} \geq \exp(P_{\text{mean}} + n\,P_{\text{std}})$ for a chosen
coverage multiplier $n$. Three values were tested, $\sigma_{\max} \in \{5,\,7,\,9\}$,
with $\sigma_{\min} = 0.002$ fixed throughout.
\newpage
% Required packages: booktabs, multirow, array
% \usepackage{booktabs, multirow, array}
\begin{table}[htbp]
\centering
\footnotesize
\renewcommand{\arraystretch}{1.35}
\setlength{\tabcolsep}{5pt}
\begin{tabular}{
  >{\itshape}p{2.0cm}
  p{3.0cm}
  p{4.5cm}
  p{4.5cm}
}
\toprule
\normalfont\textbf{Category}
  & \textbf{Parameter}
  & \textbf{Phase 1 \textemdash{} Replication}
  & \textbf{Phase 2 \textemdash{} EDM-CSDI} \\
\midrule
%--- Data & preprocessing ------------------------------------------------
\multirow{5}{2.0cm}{Data \& preprocessing}
  & Channels $K$
  & $K=1$ (close log-return)
  & $K=4$ (OHLC) \\[3pt]
  & Preprocessing
  & Raw close log-return
  & OHLC channels as log-returns relative to $C_{t-1}$;
    per-channel $z$-score normalisation \\[3pt]
  & Conditioning
  & Unconditional;\; $m^{co}=\mathbf{0}$
  & \texttt{predict\_close}: $\{O,H,L\}$ observed,
    $\{C\}$ masked \\[3pt]
  & Window length $L$
  & $2048$ (${\approx}10$ years)
  & $512$ (${\approx}2$ years) \\[3pt]
  & Stride $s$
  & $400$
  & $100$ \\
\midrule
%--- Architecture (shared) -----------------------------------------------
\multirow{5}{2.0cm}{Architecture}                      % FIX 1: 2 -> 5
  & Model channels
  & \multicolumn{2}{c}{$128$} \\[3pt]
  & Layers
  & \multicolumn{2}{c}{$6$} \\[3pt]
  & Diffusion emb.
  & \multicolumn{2}{c}{$256$} \\[3pt]              % FIX 3: added [3pt]
  & Time emb.
  & \multicolumn{2}{c}{$128$} \\[3pt]              % FIX 3: added [3pt]
  & Feature emb.
  & \multicolumn{2}{c}{$64$} \\                    % FIX 2: added \\
\midrule
%--- Noise schedule ------------------------------------------------------
\multirow{3}{2.0cm}{Noise schedule}
  & $\sigma_{\max}$ / $\sigma_{\min}$
  & VE:\;$\sigma_{\max}{=}1,\;\sigma_{\min}{=}0.01$;\;
    GBM:\;$\sigma_{\max}{=}10,\;\sigma_{\min}{=}0.1$
  & \textemdash \\[3pt]
  & $\beta_{\max}$ / $\beta_{\min}$
  & VP:\;$\beta_{\max}{=}8,\;\beta_{\min}{=}0.01$
  & \textemdash \\[3pt]
  & $P_{\text{mean}}$ / $P_{\text{std}}$
  & \textemdash
  & $P_{\text{mean}}{=}{-1.4},\;P_{\text{std}}{=}1.6$ \\
\midrule
%--- Training ------------------------------------------------------------
\multirow{4}{2.0cm}{Training}                          % FIX 4: 3 -> 4
  & LR schedule
  & \multicolumn{2}{c}{cosine annealing} \\[3pt]
  & $\eta_{\max}$ / $\eta_{\min}$
  & $5\!\times\!10^{-4}$ / $1\!\times\!10^{-6}$
  & $1\!\times\!10^{-4}$ / $1\!\times\!10^{-6}$ \\[3pt]
  & Epoch
  & $1000$
  & $500$ \\
  & Batch size
  & \multicolumn{2}{c}{$64$} \\[3pt]
  & Optimizer
  & \multicolumn{2}{c}{\texttt{AdamW}} \\
  & Early-stopping
  & \textemdash
  & validation-loss; 200 epochs\\
\midrule

\multirow{3}{2.0cm}{Sampling}
  & Method
  & \makecell[l]{Euler--Maruyama SDE\\Probability-flow ODE}
  & \makecell[l]{Heun ODE\\Stochastic Heun SDE} \\
  & $\sigma_{\max}$ / $\sigma_{\min}$
  & from \emph{Noise schedule}
  & $\sigma_{\max}=7.0$, $\sigma_{\min}=0.002$\\
\bottomrule
\end{tabular}
\caption{Training setup for Phase~1 (replication of Kim et al.\ \cite{kim2025diffusion})
and Phase~2 (EDM-CSDI). Parameters shared across both phases appear in merged cells. \emph{emb.} is short for embeddings.}
\label{tab:training_setup}
\end{table}